\section{Parallel Deciders, Protocols, and Role Cues}
\label{sec:setup}

\subsection{The Agent and Its Action Interface}
\label{sec:interface}

A tool-calling agent receives a request $q$ and a set of tool schemas $\mathcal{F}$, and acts by issuing a set of calls $A=\{a_1,\dots,a_n\}$, where each call $a_i=(f_i,\theta_i)$ names a tool $f_i\in\mathcal{F}$ and binds its arguments $\theta_i$. When $n>1$, the turn is a joint action. We score it as a set, as BFCL does~\cite{patil2025bfcl}: the turn is correct if its calls can be matched one to one with the reference calls, whatever their order.

The dLLM agent writes $A$ as a JSON array onto a canvas of mask tokens and fills the canvas over several denoising steps. To make every action well formed, its interface fixes the syntax of the turn in advance with a \emph{skeleton}: the number of calls, the tool names, the argument names, and all brackets and quotes are written onto the canvas, and only the argument values are masked. Each value occupies a \emph{slot} of $\ell$ mask tokens. The slots that hold the same argument in different calls of the same tool form a \emph{sibling group}. Our skeleton is an oracle, built from the reference calls, so that every remaining error is an error in a value; Section~\ref{sec:discussion} discusses what the oracle hides.

The skeleton writes the closing delimiter of each value after its slot: the closing quote of a string, or the comma or brace that follows any other value. A value ends early when the model writes its own closing token inside the slot; the rest of the slot is then set to padding (the model's end-of-sequence token), and both are dropped from the output. Section~\ref{sec:length} also tests a variant in which the slot holds the closing token itself and is padded with spaces. The autoregressive (AR) reference agent uses the same skeleton: it writes each value until its closing token and is cut off at the slot length.

\subsection{Deciders and Protocols}
\label{sec:protocol}

At each denoising step, the dLLM predicts every masked position from the current canvas and commits some of them. A position committed in a step is chosen from its own predicted distribution given the canvas \emph{before} that step, so positions committed in the same step cannot see each other's choices. We call the masked positions \emph{slot-level deciders}. They are positions of one network without preferences or beliefs of their own: they share one policy and differ only in their position on the canvas and in what has been committed when they move. A parallel turn is thus a common-payoff team decision problem~\cite{marschak1972teams} with a shared policy. We use the terms of game theory for its information structure, who moves when and who sees what, not for any reasoning by the deciders.

To compare, we give the same requests to \emph{teams of $n$ LLM agents}, one per call. Each agent is a separate generation of the same model, sees a skeleton of its own call only, and receives the request followed by a note on the protocol; the supplementary material gives the texts. A \emph{protocol} says which deciders move together and what each sees:
\begin{itemize}
\item \emph{Canvas.} With $k$ tokens per step in confidence order, the $k$ most confident masked positions move together. With $k=1$, one decider moves at a time and sees every earlier commitment; we call this \emph{one token per step}. Left-to-right decoding (LTR) fixes the order of moves. A confidence threshold $\tau$~\cite{wu2026fastdllm} lets every position whose top probability exceeds $\tau$ move at once, and the most confident position alone when none does. \llada decodes in blocks of 32 tokens, causal across blocks; only 3.1\% (\texttt{parallel}) and 2.3\% (\texttt{parallel\_multiple}) of the pairs of sibling slots share a block, so its calls almost always take turns.
\item \emph{Teams.} Agents act at the same time, seeing no other call, or in turns, seeing the calls made so far. They are anonymous (``You are one of $n$ assistants''), numbered (``You are assistant $i$ of $n$''), or numbered and told the rule (``By convention, assistant $i$ makes the $i$-th of the calls, in the order in which the request mentions them'').
\end{itemize}

\subsection{Role Cues and Predictions}
\label{sec:predictions}

The sibling deciders of a turn must take distinct entities, each the one its call needs. A shared policy can set them apart only through something they share that orders the candidates, or something that differs between them. We consider three such cues:
\begin{enumerate}
\item \emph{Mention order.} The request names its entities in some order. Taking the $i$-th mentioned entity for the $i$-th call is a convention that any decider who knows its index can follow, and on set-scored requests it is a focal point among equally correct assignments~\cite{schelling1960strategy}.
\item \emph{Role.} A slot's position on the shared skeleton, or an agent's number in its prompt. A role sets deciders apart only if the policy's output depends on it, and helps only if something orders the candidates: ``you are the second call'' does not say which of all US cities is the second.
\item \emph{Slot length.} Sibling values usually differ in length. In training, every mask hides exactly one token, so the mask count is a truthful signal of a value's length in the sense of Lewis~\cite{lewis1969convention}; an interface that sets slot lengths sends this signal whether or not it is true.
\end{enumerate}

\noindent These cues give three predictions, which Sections~\ref{sec:ordered}--\ref{sec:length} test.
\begin{itemize}
\item[\pred{1}] \emph{Order coordinates.} When the request orders its entities, same-step slots take the $i$-th entity for the $i$-th call at any $k$, so committing many tokens per step costs little, even when slot lengths cannot tell the calls apart.
\item[\pred{2}] \emph{Symmetric roles collide.} When nothing orders the candidates, deciders that move at the same time take the same value, and deciders that take turns do not.
\item[\pred{3}] \emph{Length is a signal.} A decider reads its mask count as the length of its value: right lengths set siblings apart, and wrong lengths give a slot a longer version of its own value or a sibling's value that fits, at any $k$, including one token per step and a single agent.
\end{itemize}
If a number in a prompt acts as a role, \pred{1} and \pred{2} also hold for numbered agents, which would take the $i$-th entity of an ordered request or of a list. Sections~\ref{sec:ordered} and~\ref{sec:symmetric} test this.

\subsection{Models, Data, and Measures}
\label{sec:models}

\paragraph{Models.} \dream~\cite{ye2025dream} (Dream-v0-Instruct-7B) is a dLLM with full bidirectional attention, initialized from \qwen-7B. \llada~\cite{bie2025llada2} (LLaDA2.0-mini, a mixture of experts with 16B parameters, 1.4B of them active per token) is a block-diffusion dLLM that supports tool calling. The AR reference is \qwen-7B-Instruct~\cite{qwen2024qwen25}, the instruction-tuned version of the base model that \dream is initialized from. All models see the same system prompt, which lists the tools as JSON and asks for a JSON array of calls. All decode greedily, except the runs at $T=0.7$ in Section~\ref{sec:symmetric}.

\paragraph{Data.} We use the \texttt{parallel} and \texttt{parallel\_multiple} categories of BFCL v4~\cite{patil2025bfcl}: 400 requests with 2 to 8 calls (2.9 on average) and 3,063 value slots. \llada's runs with wrong slot lengths use 100 of them, 50 evenly spaced in each category. In the \emph{equal-length} requests (213 for \dream's tokenizer, 216 for \llada's), each argument has slots of one length in every call that uses it. In the 104 \emph{team-symmetric} requests, the sibling calls of each tool have the same arguments and slot lengths, so neither the skeleton nor the lengths tell apart the agents that call the same tool. For requests in which nothing tells the calls apart, we add 105 \emph{choose-N} requests (Section~\ref{sec:symmetric}).

\paragraph{Slot lengths.} \emph{Exact}: each slot has as many masks as its own reference value has tokens, tokenized on its own (the first acceptable value where BFCL lists several). \emph{Surplus} $+s$: every slot gets $s\in\{1,2,4,8\}$ extra masks (\llada: no $s=4$, and $k=16$ only at $s=1$). \emph{Swap}: slot lengths are rotated by one position within each sibling group, on the requests where this changes a slot (165, and 34 for \llada). \emph{One slot longer}: in each sibling group with unequal lengths, one slot gets the length of the longest sibling value, and all others stay exact. These four derive from the reference values: they are diagnostic interventions, not lengths an agent could set. \emph{Estimate}: lengths predicted by the model itself in one forward pass (Section~\ref{sec:design}).

\paragraph{Measures.} \emph{Set accuracy} is the BFCL verdict. The \emph{cross-call error rate} (\ccer) is the share of requests with at least one cross-call error, found by an optimal one-to-one matching of predicted and reference calls: a duplicated or missing call, a value that belongs to another call (cross-binding), a string spliced from two calls' values (chimera), or an argument that the reference shares across calls but the prediction does not. At the slot level, we read each slot's value from its own tokens and classify it as the slot's own reference value, a sibling's value (\emph{fit} if it has exactly the slot's length), an \emph{overfill} (the own value followed by more), a truncation of the own value, or anything else. For paired comparisons on the same requests we give the difference with a 95\% bootstrap interval over requests and an exact McNemar test, and for single rates a Wilson interval; tests are not corrected for multiplicity.
